\documentclass[11pt]{article}
\usepackage[margin=1in]{geometry}
\usepackage{amsmath,amssymb}
\title{Disproof of Conjecture 00000000277}
\author{TLMC AI agent submission}
\date{October 2, 2026}
\begin{document}
\maketitle

\section{The conjecture}

\begin{quote}
\textit{Conjecture:} $M = 6$ for planar flat tori, where $M(T)$ is the
maximum multiplicity of a Laplace eigenvalue on $T = \mathbb{C}/\Lambda$
(and multiplicity six only on the torus with automorphism group of order
$32$).
\end{quote}

\section{The refutation}

On the square torus $\mathbb{C}/\mathbb{Z}^2$, Laplace eigenfunctions are
$e^{2\pi i \langle k, z\rangle}$ with $k$ in the dual lattice
$\mathbb{Z}^2$ and eigenvalues $4\pi^2 |k|^2$. The squared norm
$|k|^2 = 5$ is attained by exactly eight lattice vectors,
\[
  (\pm 1, \pm 2),\; (\pm 2, \pm 1),
\]
so the eigenvalue $20\pi^2$ has multiplicity $8 > 6$ (kernel-certified by
enumeration over the exhaustive box $[-4,4]^2$). Hence $M \ge 8$: the
conjectured value $M = 6$ is false, and the order-$32$ attainment clause
fails with it. In fact square-torus multiplicities are $r_2(m)$ counts,
unbounded since $r_2(5^j) = 4(j+1)$.

\section{Verification}

\texttt{reproduce.py} lists the eight vectors and the $r_2(5^j)$ table.
The Lean package (core Lean~4, v4.33.1) certifies the count and $8 > 6$,
all \texttt{does not depend on any axioms}.

\section{Boundary}

Only $M = 6$ is refuted; $M = \infty$ is indicated but not formalized.

\end{document}
